\section{Rings and Modules of Fractions}

\begin{exercise}
	Element-wise. Only if part requires the finitely generated property.
\end{exercise}

\begin{exercise}
	Element-wise. Hint.
\end{exercise}

\begin{exercise}\label{ex3.3}
	Element-wise.
\end{exercise}

\begin{exercise}
	Element-wise.
\end{exercise}

\begin{exercise}
	Element-wise. Counter example is Boolean ring.
\end{exercise}

\begin{exercise}\label{ex3.6}
	Use the fact that for any multiplicatively closed set $S$, if $A-S$ is an ideal, it's automatically prime.
\end{exercise}

\begin{exercise}\label{ex3.7}
	\begin{enumerate}
		\item[(i),(ii)]More specifically, $A-S=\bigcup_{\mathfrak{p}\cap S=\varnothing}\mathfrak{p}$.
		\item[(iii)]$\mathfrak{p}\cap(1+\mathfrak{a})=\varnothing$ $\Leftrightarrow$ $\mathfrak{a}+\mathfrak{p}\neq(1)$ $\Leftrightarrow$ $\mathfrak{a}+\mathfrak{p}\subseteq\mathfrak{m}$ for some maximal ideal $\mathfrak{m}$. Note that the union of a chain of prime ideals is only a maximum ideal.
	\end{enumerate}
\end{exercise}

\begin{exercise}
	Element-wise. (v) $\Rightarrow$ (ii): If $t/1$ is not a unit in $S^{-1}A$, it contains in a prime ideal $\mathfrak{q}\subseteq S^{-1}A$, so $t\in T\cap\mathfrak{q}^c$.
\end{exercise}

\begin{exercise}
	Use \Cref{ex3.6}. Element-wise.
\end{exercise}

\begin{exercise}\label{ex3.10}
	(i) \Cref{ex2.27} says for any element $x\in A$, $x=ax^2$ for some $a\in A$. (ii) $\Rightarrow$: Use \Cref{ex2.28}. $\Leftarrow$: Consider an $A_\mathfrak{m}$-module. Use \Cref{3.10}.
\end{exercise}

\begin{exercise}\label{ex3.11}
	Let $A$ be a ring. Then the following are equivalent:
	\begin{enumerate}[label=\textup{(\roman*)}]
		\item $A/\mathfrak{R}$ is absolutely flat.
		\item Every prime ideal of $A$ is maximal.
		\addtocounter{enumi}{1}
		\item $\operatorname{Spec}(A)$ is Hausdorff.
	\end{enumerate}
	(i) $\Rightarrow$ (ii): Use \Cref{ex2.27} on $A/\mathfrak{p}$. (ii) $\Rightarrow$ (i): Proof that $\left(A/\mathfrak{R}\right)_{\overline{\mathfrak{p}}}$ has no nilpotent, then use \Cref{ex1.10,ex3.10}. (iv) $\Rightarrow$ (ii): There exist $f_1\in\mathfrak{p}_1$ and $f_2\in\mathfrak{p}_2$ such that $f_1+f_2=1$, so $f_2\in X_{f_1}$ and $f_1\in X_{f_2}$. \textcolor{magenta}{need to finish the proof}
\end{exercise}

\begin{exercise}\label{ex3.12}
	Element-wise. (iv) Use \Cref{3.5}.
\end{exercise}

\begin{exercise}\label{ex3.13}
	Element-wise.
\end{exercise}

\begin{exercise}\label{ex3.14}
	Hint. Use \Cref{3.8}.
\end{exercise}

\begin{exercise}\label{ex3.15}
	Hint.
\end{exercise}

\begin{exercise}\label{ex3.16}
	(ii) $\Rightarrow$ (iii): \Cref{1.17}. (iii) $\Rightarrow$ (iv) $\Rightarrow$ (v): Use hint. (v). Take $M=A/\mathfrak{a}$. Then $M_B=B/\mathfrak{a}^e$.
\end{exercise}

\begin{exercise}\label{ex3.17}
	
\end{exercise}

\begin{exercise}\label{ex3.18}
	Use \Cref{ex3.3} to proof $B_\mathfrak{q}\cong\left(B_\mathfrak{p}\right)_{\mathfrak{q}B_\mathfrak{p}}$. By \Cref{3.3} and \Cref{3.5}, $B_\mathfrak{q}$ is flat over $B_\mathfrak{p}$. 
\end{exercise}

\begin{exercise}\label{ex3.19}
	\begin{enumerate}[label=\textup{(\roman*)},ref=\theexercise(\roman*)]\crefalias{enumi}{exercise}
		\addtocounter{enumi}{4}
		\item\label{ex3.19.5} $M\cong\bigoplus_{i=1}^nA/\mathfrak{a}_i$, so $\operatorname{Supp}(M)=\bigcup_{i=1}^n\operatorname{Supp}(A/\mathfrak{a}_i)=\bigcup_{i=1}^nV(\mathfrak{a}_i)=V\left(\bigcap_{i=1}^n\mathfrak{a}_i\right)$. Proof that $\bigcap_{i=1}^n\mathfrak{a}_i=\operatorname{Ann}(M)$.
		\addtocounter{enumi}{1}
		\item Use \Cref{ex2.2}.
		\item
	\end{enumerate}
\end{exercise}

\begin{exercise}
	Counter example is $A[x]/\left(x^2\right)$.
\end{exercise}

\begin{exercise}\label{ex3.21}
	(i)-(iii) Use one-to-one order preserving correspondences. (iv) $A_\mathfrak{p}/\mathfrak{p}A_\mathfrak{p}=(A/\mathfrak{p})_\mathfrak{p}$.
\end{exercise}

\begin{exercise}
	Element-wise.
\end{exercise}

\begin{exercise}
	(i) Use \Cref{ex1.17}. (v) Define two homomorphisms by the universal properties of direct limit and local ring.
\end{exercise}

\begin{exercise}
	
\end{exercise}
